% ID: FIS-TER-GAS-004
% Titolo: Calore e lavoro in una trasformazione ciclica
% Disciplina: Fisica
% Area: Termologia
% Argomento: Gas perfetti
% Sottoargomento: Primo principio in un ciclo termodinamico
% Classe: Quarta liceo scientifico
% Difficolta: 3
% Tipo: quesito_teorico
% Risultato: in un ciclo Delta U = 0, quindi Q = L
% Tempo_stimato: 6 min
% Competenze: interpretazione del primo principio, argomentazione, riconoscimento di un ciclo
% Prerequisiti: primo principio della termodinamica, energia interna, lavoro
% Tag: termologia, gas_perfetti, ciclo_termodinamico, calore, lavoro
% Fonte: verifica_fisica_gas_perfetti_anonimizzata
% Autore: Riccardo Carli
% Licenza: CC-BY-SA-4.0

\begin{esercizio}
In una trasformazione ciclica, il calore scambiato dal sistema e uguale al
lavoro fatto dal sistema? Motiva la tua risposta usando il primo principio
della termodinamica.
\end{esercizio}

\begin{soluzione}
Si, se si usa la convenzione in cui \(L\) e il lavoro fatto dal sistema. Il
primo principio della termodinamica si scrive:
\[
\Delta U=Q-L.
\]

In una trasformazione ciclica lo stato finale coincide con quello iniziale.
L'energia interna e una funzione di stato, quindi dopo un ciclo completo:
\[
\Delta U=0.
\]

Allora:
\[
0=Q-L
\qquad\Rightarrow\qquad
Q=L.
\]

Questo significa che il calore netto assorbito dal sistema in un ciclo e uguale
al lavoro netto compiuto dal sistema. Se si usa una convenzione diversa per il
segno del lavoro, la relazione cambia forma ma il contenuto fisico resta lo
stesso.
\end{soluzione}
